
Code generation models are evaluated on multiple benchmarks under the assumption that different benchmarks measure distinct competencies. We empirically test this assumption by analyzing execution trace features across model populations on HumanEval and MBPP benchmarks. We extracted standardized features (pass@k metrics, runtime quartiles, error distributions) for 8 code generation models with 100\% completeness across 14 model-benchmark pairs. Analysis reveals perfect ranking correlation (Spearman ρ = 1.000, p < 0.0001, n=6 models) between HumanEval and MBPP despite substantial distributional divergence (KL divergence = 18.4). This finding indicates that the two benchmarks, while differing in statistical properties, produce identical model competency orderings. The result refutes the hypothesis that benchmark design differences (algorithmic clarity vs. practical patterns) create distinctive evaluation signatures. We interpret this as evidence that execution-based benchmarks may measure a unidimensional competency construct with different difficulty calibrations. The small sample size (n=6) limits generalizability, and the analysis is restricted to two benchmarks. Results inform benchmark selection for model evaluation and suggest that evaluation diversity may require task-type variety (generation, understanding, repair) rather than multiple instances of the same task type.
